\documentclass[10pt,a4paper]{amsart}
\usepackage[margin=19mm]{geometry}
\usepackage{amsmath,amssymb,mathtools}
\usepackage[scr=rsfs,cal=euler]{mathalfa}
\usepackage{xfrac}
\usepackage[unicode,hidelinks,bookmarks=false]{hyperref}
\newcommand{\ie}{i.e.\ }
\newcommand{\cf}{cf.\ }
\newcommand{\resp}{respectively\ }
\newcommand{\Subsection}{\subsection}
\providecommand{\nobreakdash}{\nobreak\hbox{-}}
\setlength{\emergencystretch}{2em}
\multlinegap0pt
\usepackage{tikz}
\usetikzlibrary{cd}
\usepackage{enumerate}
\usepackage{amssymb}
\usepackage{tikz}
\usepackage{tcolorbox}
\usepackage{xcolor}
\usepackage{cancel}
\newtheorem{definition}{Definition}
\newtheorem{example}{Example}
\newtheorem{proposition}{Proposition}
\usepackage{cleveref}
\crefname{definition}{Definition}{Definitions}
\crefname{example}{Example}{Examples}
\crefname{proposition}{Proposition}{Propositions}
\newcommand{\fX}{{\mathfrak X}}
\DeclareMathOperator{\id}{id}
\newcommand{\trans}[1]{{\lfloor #1 \rfloor}}
\title{The \'{e}tale Brauer-Manin obstruction for classifying stacks}
\author{Ajneet Dhillon, Nicole Lemire, Jonathan Martin, Yidi Wang}
\date{Source: arXiv:2512.01014v2 (CC BY 4.0)}
\begin{document}
\maketitle

\noindent\emph{Source and licence.} The \'{e}tale Brauer-Manin obstruction for classifying stacks. Version arXiv:2512.01014v2 (CC BY 4.0), distributed by arXiv under the Creative Commons Attribution 4.0 International licence, http://creativecommons.org/licenses/by/4.0/, as recorded in the arXiv licence metadata for this exact version and shown on the abstract page https://arxiv.org/abs/2512.01014v2. Full licence notice: \texttt{licenses/arxiv-2512.01014-CC-BY-4.0.txt}.\par\smallskip

\noindent\textit{Editorial scope.} Counted unit: the article's proposition p:innerAction (source0.tex lines 1181--1195). The article's complete original proof of it is delivered with it (source0.tex lines 1198--1345, one \texttt{proof} environment of 148 lines). No mathematics is written, repaired or re-derived: the statement and the proof are the article's own lines, in the article's own notation, and no retained line is substituted or re-flowed.  Retained uncounted as the source-local context the proof needs: lines 939--971; lines 1126--1160.  Nothing was omitted from any retained span, so the package carries no omission record.  No retained line contains a citation, so the wrapper carries no bibliography and no bibliography entry.  The retained lines contain the article's own diagram code, which the wrapper typesets with the standard tikz packages; no image file is delivered and the package declares no asset dependency.  Editorial formatting only, in addition: an AMS \texttt{amsart} preamble that restates the article's own declarations and macros used by the retained lines, namely the environments \texttt{definition}, \texttt{example}, \texttt{proposition} on separate counters, together with the standard packages those lines need.  The retained environments print in this document as Definition 1, Example 1, Proposition 1 in order of appearance, whereas the article prints the same items with the numbering of its own full document.  The complete native source of the article as distributed in the arXiv e-print for this exact version is bundled unchanged as \texttt{sources/190097/source0.tex}.
\begin{definition}\label{d:weakAction}    
    Let $\fX$ be a category fibered in groupoids. 
    A right action of $G$ on $\fX$ is a triple 
    $(\mu, \alpha, a)$ where $\mu: \fX\times G\to \fX$ is a morphism such that 
    the following diagrams are 2-commutive:
    $$
    \begin{tikzcd}
        {\fX\times G \times G} & {\fX\times G} && {\fX\times G} & \fX \\
        {\fX\times G} & \fX && \fX
        \arrow["{\id_\fX\times m}", from=1-1, to=1-2]
        \arrow["{\mu\times \id_G}"', from=1-1, to=2-1]
        \arrow["\mu", from=1-2, to=2-2]
        \arrow["\mu", from=1-4, to=1-5]
        \arrow["\alpha", Rightarrow, from=2-1, to=1-2]
        \arrow["\mu"', from=2-1, to=2-2]
        \arrow["{\id_\fX\times \{e\}}", from=2-4, to=1-4]
        \arrow[""{name=0, anchor=center, inner sep=0}, "{\id_\fX}"', from=2-4, to=1-5]
        \arrow["a"', Rightarrow, from=0, to=1-4]
    \end{tikzcd}
    $$
    where $e$ is the identity of $G$. 
    In what follows we will write $\mu(x,g)=x\cdot g$ where $x$ is an object (or a morphism) 
    of  $\fX$. Given $(x,g,h)\in \fX\times G\times G$ we write 
    $\alpha^{x}_{g,h}:x\cdot g \cdot h \to x \cdot gh$ and $a^x:x \to x\cdot e$
    for the value of the natural transformations $\alpha$ and $a$ respectively on 
    these objects. The above data is subject to the axioms:
    \begin{enumerate}
        \item $\alpha_{gh,k}^{x}\circ (\alpha_{g,h}^{x}\cdot k) = \alpha_{g,hk}^{x} \circ \alpha_{h,k}^{x\cdot g}$ as morphisms $x\cdot g \cdot h \cdot k \to x\cdot ghk$.
        \item $1_{x\cdot g}=(\alpha^{x}_{e,g})\circ (a^{x}\cdot g)$ as morphisms $x\cdot g \to x \cdot g$.
        \item $1_{x\cdot g}=(\alpha^{x}_{g,e})\circ (a^{x\cdot g})$ as morphisms $x\cdot g \to x \cdot g$. 
    \end{enumerate}
    A \emph{strict action of $G$} on $\fX$ is a weak action where $\alpha$ and $a$ are both the identity. 
\end{definition}

\begin{example}\label{e:inner} This example is well-known. We present it here in terms of the sheafification of transformation groupoids. 
    Let $G$ be a linear algebraic group over a field $k$ and $P$ a left $G$-torsor over ~$k$. 
    Then there is a right action of $G$ on $G \times P$ given by:
    $$
    (g,p) \cdot h := (g^h,h^{-1}p) \quad \text{ where } g^{h}=h^{-1}gh 
    $$
    The space obtained by taking the quotient of this action is called the \emph{inner form}
    associated to $P$, and is denoted by $G^{P}$. It is in fact a group. To construct the 
    multiplication map consider the morphism 
    $$
    G\times P \times G\times P \to G\times P
    $$
    which, for a $k$-scheme $S$, is given on $S$-points by:
    $$
    (g_{1},p,g_{2}, p') \mapsto (g_{1} g_{3}^{-1}g_{2}g_3,p)
    $$ where $g_3$ is the unique $S$-point of $G$ such that $p' = g_3\cdot p$. 

    There is an action of $G \times G$ on $G\times P\times G\times P$ and an action of 
    $G$ on $G\times P$. One then checks that this morphism is $\pi_{2}$-equivariant, where 
    $\pi_2:G\times G\to G$ is the projection onto the second factor. 

    This gives a morphism of transformation groupoids 
    $$
    \trans{G\times P\times G \times P/G^2}\to \trans{G\times P/G}
    $$
    and hence a morphism $G^{P}\times G^{P}\to G^{P}$ after sheafifying. 
    Similarly, to construct inversion and unit maps, one considers the morphisms 
    $\iota:G \times P \to G \times P$ given by $(g,p) \mapsto (g^{-1},p)$, and 
    $\eta:P \to G \times P$ given by $p \mapsto (e,p)$ respectively. Since these
    morphisms are $G$-equivariant when $P$ is given a right $G$-action by 
    inverting its left $G$-action, it follows that these morphisms also give 
    morphisms of transformation groupoids. After sheafifying, one gets 
    morphisms $G^P \to G^P$ and $* \to G^P$. One then readily checks that 
    $G^{P}$ becomes a group under these morphisms.
\end{example}

\begin{proposition}\label{p:innerAction}
    In the situation of the above examples,
    the morphism 
    $$
    \rho :  \fX\times P \times  G\times P\to \fX\times P
    $$
    given by 
    $$
    (x,p,g,p')\mapsto (x\cdot (h^{-1}gh),p), 
    $$ where $h$ is the unique point of $G$ that satisfies $h\cdot p = p'$,
    induces a weak group action 
    $$
    \bar{\rho}: \fX^{P}\times G^{P}\to \fX^{P}.
    $$
\end{proposition}

\begin{proof}
    There is a $G\times G$ action on $\fX\times P \times  G\times P$ given by combining the actions of the previous two examples. Explicitly
    $$
    (x,p, g, hp)\cdot (s_{1},s_{2}) = (x\cdot s_{1}, s_{1}^{-1}p,g^{s_{2}}, s_{2}^{-1}hp). 
    $$
    We need to show that $\rho$ is $\pi_1$-equivariant where $\pi_1:G\times G\to G$ is the 
    projection onto the first factor. That is , if we let $(\mu_{1},\alpha,a)$ 
    (resp. $\mu_{2}$) be the action of $G$ on $\fX$ (resp. $P$) then we need to give a 
    2-morphism $\sigma$ such that the following diagram 2-commutes:
    $$
    \begin{tikzcd}
        {\fX\times P\times G \times P \times G \times G} & {\fX\times P \times G \times P} \\
        {\fX\times P \times G} & {\fX\times P}
        \arrow[from=1-1, to=1-2]
        \arrow["{\rho\times\pi_1}"', from=1-1, to=2-1]
        \arrow["\rho", from=1-2, to=2-2]
        \arrow["\sigma", Rightarrow, from=2-1, to=1-2]
        \arrow["\mu"', from=2-1, to=2-2]
    \end{tikzcd}
    $$
    Given $(x,p,g,hp,s_1,s_2) \in \fX\times P \times G \times P \times G \times G$,
    then by going along the bottom path of the above diagram, we get 
    $(x\cdot g^{h} \cdot s_{1},s_{1}^{-1}p)$. Going along the top we get
    $(x\cdot s_{1}\cdot g^{hs_1}, s_{1}^{-1}p)$. Since the $P$ component of
    both of these pairs are the same, to give an isomorphism between them it suffices 
    to give an isomorphism 
    $$x\cdot g^{h}\cdot s_{1} \to x\cdot s_{1}\cdot g^{hs_1}$$
    in $\fX$. One then checks that if 
    $\omega_{s}^{x,g} = (\alpha_{s,g^{s}}^{x})^{-1} \circ \alpha_{g,s}^{x}$
    then $\omega_{s_{1}}^{x,g^{h}}$ is the desired isomorphism. Thus 
    $$\sigma_{(s_1,s_2)}^{(x,p,g,hp)} = \omega_{s_{1}}^{x,g^{h}} \times \id_{s_{1}^{-1}p}$$
    
    This proves the existence of the morphism $\bar{\rho}$. It remains to check that 
    $\bar{\rho}$ can be extended to an action $(\bar{\rho}, \bar{\lambda}, \bar{l})$. 
    To do this, we will construct the appropriate natural isomorphisms on 
    $\fX \times P \times G\times P$ and check that they are equivariant. Let $m$
    denote the morphism $G \times P \times G \times P \to G \times P$
    constructed in \Cref{e:inner} which descends to multiplication 
    $G^{P} \times G^{P} \to G^{P}$.

    To construct $\lambda$ is to give a 2-morphism such that the following 
    diagram 2-commutes:
    $$
    \begin{tikzcd}
        {\fX\times P\times G \times P \times G \times P} & {\fX\times P \times G \times P} \\
        {\fX\times P \times G \times P} & {\fX\times P}
        \arrow["{\id_{\fX\times P}\times m}", from=1-1, to=1-2]
        \arrow["{\rho\times\id_{G\times P}}"', from=1-1, to=2-1]
        \arrow["\rho", from=1-2, to=2-2]
        \arrow["\lambda", Rightarrow, from=2-1, to=1-2]
        \arrow["\rho"', from=2-1, to=2-2]
    \end{tikzcd}
    $$
    Given 
    $(x,p,g_{1},h_{1}p,g_{2},h_{2}p) \in \fX\times P\times G \times P \times G \times P$,
    then going along the top path of the above diagram we get
    $(x\cdot g_{1}^{h_{1}}g_{2}^{h_{2}},p)$. Going along the bottom 
    path we get $(x \cdot g_{1}^{h_{1}}\cdot g_{2}^{h_{2}}, p)$.
    Similar to defining $\sigma$, to give an isomorphism from the latter to the former
    it suffices to give an isomorphism
    $$x \cdot g_{1}^{h_{1}}\cdot g_{2}^{h_{2}} \to x\cdot g_{1}^{h_{1}}g_{2}^{h_{2}} $$
    in $\fX$. By definition, 
    $\alpha_{g_{1}^{h_{1}}, g_{2}^{h_{2}}}^{x}$
    is the desired isomorphism. So 
    $$\lambda^{(x,p,g_1,h_1p,g_2,h_2p)}= \alpha_{g_{1}^{h_{1}},g_{2}^{h_{2}}}^{x} \times \id_{p}$$

    To show that $\lambda$ descends as desired requires us to check that it is 
    $\pi_1$-equivariant, where $\pi_1:G^{3} \to G$ is the first projection,
    and $\fX\times P \times G \times P \times G \times P$ is given an action by 
    $G^{3}$ action by combining the actions of $G$ on $\fX \times P$ and 
    $G \times P$. If one unpacks all of the definitions, constructions and computations, 
    this reduces to checking that for all 
    $(x,p,g_{1},h_{1}p,g_{2},h_{2}p,s) \in \fX\times P \times G \times P \times G \times P \times G$
    the following diagram commutes in $\fX$:
    $$
    \begin{tikzcd}[column sep=huge]
        {x\cdot g_{1}^{h_{1}}g_{2}^{h_{2}}\cdot s} & {x\cdot s\cdot g_{1}^{h_{1}s}g_{2}^{h_{2}s}} \\
        {x\cdot g_{1}^{h_{1}}\cdot g_{2}^{h_{2}}\cdot s} & {x\cdot s\cdot g_{1}^{h_{1}s}\cdot g_{2}^{h_{2}s}}
        \arrow["{\omega_{s}^{x,g_{1}^{h_{1}}g_{2}^{h_{2}}}}", from=1-1, to=1-2]
        \arrow["{(\alpha_{g_{1}^{h_{1}},g_{2}^{h_{2}}}^{x})\cdot s}", from=2-1, to=1-1]
        \arrow["{(\omega_{s}^{x,g_{1}^{h_{1}}}\cdot g_{2}^{h_{2}s})\circ\omega_{s}^{x\cdot g_{1}^{h_{1}},g_{2}^{h_{2}}}}"', from=2-1, to=2-2]
        \arrow["{\alpha_{g_{1}^{h_{1}s},g_{2}^{h_{2}s}}^{x\cdot s}}"', from=2-2, to=1-2]
    \end{tikzcd}
    $$
    Recall $\omega_{s}^{x,g} = (\alpha_{s,g^{s}}^{x})^{-1} \circ \alpha_{g,s}^{x}$.
    The commutativity of this diagram can then be shown by repeated uses of (1) in 
    \Cref{d:weakAction}.
    \iffalse
    \mac{The following is the computations needed to show the diagram commutes. This should not be in the 
    final version, but I will leave it here in the comments for reference}
    Going along the top we get:
    $$(\alpha_{s,g_{1}^{h_{1}s}g_{2}^{h_{2}s}}^{x})^{-1}\circ (\alpha_{g_{1}^{h_{1}}g_{2}^{h_{2}},s}^{x})\circ(\alpha_{g_{1}^{h_{1}},g_{2}^{h_{2}}}^{x}\cdot s) $$
    Going along the bottom we get:
    $$(\alpha_{g_{1}^{h_{1}s},g_{2}^{h_{2}s}}^{x\cdot s}) \circ (\alpha_{s,g_{1}^{h_{1}s}}^{x}\cdot g_{2}^{h_{2}s})^{-1}\circ (\alpha_{g_{1}^{h_{1}},s}^{x}\cdot g_{2}^{h_{2}s})\circ (\alpha_{s,g_{2}^{h_{2}s}}^{x\cdot g_{1}^{h_{1}}})^{-1}\circ(\alpha_{g_{2}^{h_{2}},s}^{x\cdot g_{1}^{h_{1}}}) $$

    Applying (1) from \Cref{d:weakAction}, we have the following relations:
    $$\alpha_{g_{1}^{h_{1}},g_{2}^{h_{2}}}^{x} \cdot s = (\alpha_{g_{1}^{h_{1}}g_{2}^{h_{2}},s}^{x})^{-1}\circ (\alpha_{g_{1}^{h_{1}},sg_{2}^{h_{2}s}}^{x})\circ (\alpha_{g_{2}^{h_{2}},s}^{x\cdot g_{1}^{h_{1}}})$$
    $$\alpha_{g_{1}^{h_{1}},s}^{x} \cdot g_{2}^{h_{2}s} = (\alpha_{sg_{1}^{h_{1}s},g_{2}^{h_{2}s}}^{x})^{-1}\circ (\alpha_{g_{1}^{h_{1}},sg_{2}^{h_{2}s}}^{x})\circ (\alpha_{s,g_{2}^{h_{2}s}}^{x\cdot g_{1}^{h_{1}}})$$
    $$(\alpha_{s,g_{1}^{h_{1}s}}\cdot g_{2}^{h_{2}s})^{-1} = (\alpha_{g_{1}^{h_{1}s},g_{2}^{h_{2}s}}^{x\cdot s})^{-1}\circ(\alpha_{s,g_{1}^{h_{1}s}g_{2}^{h_{2}s}}^{x})^{-1}\circ(\alpha_{sg_{1}^{h_{1}s},g_{2}^{h_{2}s}}^{x})$$

    So, starting from the bottom path we have:
    \begin{align*}
         &(\alpha_{g_{1}^{h_{1}s},g_{2}^{h_{2}s}}^{x\cdot s}) \circ (\alpha_{s,g_{1}^{h_{1}s}}^{x}\cdot g_{2}^{h_{2}s})^{-1}\circ (\alpha_{g_{1}^{h_{1}},s}^{x} \cdot g_{2}^{h_{2}s})\circ (\alpha_{s,g_{2}^{h_{2}s}}^{x\cdot g_{1}^{h_{1}}})^{-1}\circ(\alpha_{g_{2}^{h_{2}},s}^{x\cdot g_{1}^{h_{1}}}) \\
        =&(\alpha_{g_{1}^{h_{1}s},g_{2}^{h_{2}s}}^{x\cdot s}) \circ (\alpha_{s,g_{1}^{h_{1}s}}^{x}\cdot g_{2}^{h_{2}s})^{-1}\circ (\alpha_{sg_{1}^{h_{1}s},g_{2}^{h_{2}s}}^{x})^{-1}\circ (\alpha_{g_{1}^{h_{1}},sg_{2}^{h_{2}s}}^{x})\circ (\alpha_{s,g_{2}^{h_{2}s}}^{x\cdot g_{1}^{h_{1}}})\circ (\alpha_{s,g_{2}^{h_{2}s}}^{x\cdot g_{1}^{h_{1}}})^{-1}\circ(\alpha_{g_{2}^{h_{2}},s}^{x\cdot g_{1}^{h_{1}}}) \\
        =&(\alpha_{g_{1}^{h_{1}s},g_{2}^{h_{2}s}}^{x\cdot s}) \circ (\alpha_{s,g_{1}^{h_{1}s}}^{x}\cdot g_{2}^{h_{2}s})^{-1}\circ (\alpha_{sg_{1}^{h_{1}s},g_{2}^{h_{2}s}}^{x})^{-1}\circ (\alpha_{g_{1}^{h_{1}},sg_{2}^{h_{2}s}}^{x})\circ(\alpha_{g_{2}^{h_{2}},s}^{x\cdot g_{1}^{h_{1}}}) \\
        =&(\alpha_{g_{1}^{h_{1}s},g_{2}^{h_{2}s}}^{x\cdot s}) \circ (\alpha_{g_{1}^{h_{1}s},g_{2}^{h_{2}s}}^{x\cdot s})^{-1}\circ (\alpha_{s,g_{1}^{h_{1}s}g_{2}^{h_{2}s}}^{x})^{-1}\circ(\alpha_{sg_{1}^{h_{1}s},g_{2}^{h_{2}s}}^{x})\circ (\alpha_{sg_{1}^{h_{1}s},g_{2}^{h_{2}s}}^{x})^{-1}\circ (\alpha_{g_{1}^{h_{1}},sg_{2}^{h_{2}s}}^{x})\circ(\alpha_{g_{2}^{h_{2}},s}^{x\cdot g_{1}^{h_{1}}}) \\
        =&(\alpha_{s,g_{1}^{h_{1}s}g_{2}^{h_{2}s}}^{x})^{-1} \circ (\alpha_{g_{1}^{h_{1}},sg_{2}^{h_{2}s}}^{x}) \circ(\alpha_{g_{2}^{h_{2}},s}^{x\cdot g_{1}^{h_{1}}}) \\
        =&(\alpha_{s,g_{1}^{h_{1}s}g_{2}^{h_{2}s}}^{x})^{-1} \circ (\alpha_{g_{1}^{h_{1}}g_{2}^{h_{2}},s}^{x}) \circ (\alpha_{g_{1}^{h_{1}}g_{2}^{h_{2}},s}^{x})^{-1} \circ (\alpha_{g_{1}^{h_{1}},sg_{2}^{h_{2}s}}^{x})\circ(\alpha_{g_{2}^{h_{2}},s}^{x\cdot g_{1}^{h_{1}}}) \\
        =&(\alpha_{s,g_{1}^{h_{1}s}g_{2}^{h_{2}s}}^{x})^{-1} \circ (\alpha_{g_{1}^{h_{1}}g_{2}^{h_{2}},s}^{x}) \circ (\alpha_{g_{1}^{h_{1}},g_{2}^{h_{2}}}^{x}\cdot s)
    \end{align*}
    Which is exactly the top path, so the diagram commutes as desired.
    \fi

    To construct $l$ is to give a 2-morphism such that the following diagram 2-commutes:
    $$
    \begin{tikzcd}
        {\fX\times P \times G \times P} & {\fX \times P} \\
        {\fX \times P}
        \arrow["\rho", from=1-1, to=1-2]
        \arrow["\id_{\fX \times P} \times \eta", from=2-1, to=1-1]
        \arrow[""{name=0, anchor=center, inner sep=0}, "{\id_{\fX\times P}}"', from=2-1, to=1-2]
        \arrow["l"{description}, Rightarrow, from=0, to=1-1]
    \end{tikzcd} 
    $$ Here $\eta$ is the unit map as defined in \Cref{e:inner}.
    Given $(x,p) \in \fX \times P$, then going along the top path we get $(x\cdot e,p)$ and 
    going along the bottom path we get $(x,p)$. So we use:
    $l^{(x,p)} = a^{x} \times \id_p$ to make the diagram 2-commute.
    To show that $l$ descends as desired requires us to check that it is $G$-equivariant.
    Similar to checking that $\lambda$ descends, after unpacking, this problem reduces to checking
    that for all $(x,p,g) \in \fX \times P \times G$ the following diagram commutes in $\fX$:
    $$
    \begin{tikzcd}
        {x\cdot e \cdot g} & {x\cdot g \cdot e} \\
        {x\cdot g} & {x\cdot g}
        \arrow["{\omega_{g}^{x,e}}", from=1-1, to=1-2]
        \arrow["{a^{x}\cdot g}", from=2-1, to=1-1]
        \arrow["{\id_{x\cdot g}}"', from=2-1, to=2-2]
        \arrow["{a^{x\cdot g}}"', from=2-2, to=1-2]
    \end{tikzcd}
    $$
    Since $\omega_{g}^{x,e} = (\alpha_{g,e}^{x})^{-1}\circ\alpha_{e,g}^{x}$
    this follows immediately from (2) and (3) in \Cref{d:weakAction}.

    Thus, both $\lambda$ and $l$ descend as desired. This produces the data required to 
    define a weak group action of $G^P$ on the transformation groupoid $\lfloor \fX/G\rfloor$. 
    One readily checks that it satisfies the axioms of \ref{d:weakAction}. The 
    action on quotient stacks is then obtained by stackifying.
\end{proof}

\end{document}
